\scanpage{655}
\mspage{630}

Let us reduce modulo $\mathcal{G}^{\unsure{\natural|\sharp}}$,
the relation $\alpha''\in\mathcal{G}^{\unsure{\natural|\sharp}}_\tau$
reads
\[
\dot{\alpha}''\in\{1,\dot{\tau}\}\quad\text{i.e.}\quad
\dot{\varepsilon}\,u'(\sigma)^{\cdot}\dot{\alpha}'\dot{\sigma}\dot{\tau}
\in\{1,\dot{\tau}\}
\quad\text{i.e.}\quad
\dot{\varepsilon}\,u'(\sigma)^{\cdot}\dot{\alpha}'\dot{\sigma}
\in\{1,\dot{\tau}\}.
\]
\deleted{\unsure{we would have}, since $\sigma\tau=-\tau\sigma$}
\inserted{We distinguish \unsure{again} two cases}

$1^\circ)$ \quad $\varepsilon'_3=1$ i.e. $\dot{\alpha}'=1$
\deleted{$\unsure{\dot{\varepsilon}\,u'(\sigma)^{\cdot}
\dot{\alpha}'\dot{\tau}}$}
the relation reads
\[
(*)\qquad \dot{\varepsilon}\,u'(\sigma)^{\cdot}\dot{\sigma}=1
\]
$2^\circ)$ \quad $\varepsilon'_3=-1$ i.e. $\dot{\alpha}'=\dot{\tau}$,
since $\tau\sigma=-\sigma\tau$, the relation reads
\[
(**)\qquad -\dot{\varepsilon}\,u'(\sigma)^{\cdot}\dot{\sigma}=1,
\]
thus in all cases, the relation reads
\[
(9)\qquad
\left\{
\begin{array}{l}
u'(\sigma)^{\cdot}=(\varepsilon\varepsilon'_3)^{\cdot}
\dot{\sigma}^{-1}
\quad\deleted{\unsure{\text{or}}}\ \text{i.e. again}\\
u'(\sigma)^{\cdot}=(-\varepsilon\varepsilon'_3)^{\cdot}\dot{\sigma}
\end{array}
\right.
\]
Now \deleted{let us use hypothesis b) on $u'$,}
\inserted{suppose now that we have b) for $u'$,}
which \deleted{\unsure{means}}\inserted{implies}
\[
u'(\sigma)^{\cdot}
=\varepsilon'_2\,\deleted{\text{\illegible{1}}}\dot{\sigma},
\]
we see therefore that it suffices to choose $\varepsilon$ so that
$-\varepsilon\varepsilon'_3=\varepsilon'_2$
i.e.
\[
\varepsilon=-\varepsilon'_2\varepsilon'_3
\]
thus we will have
\[
\deleted{\unsure{(10)}}\qquad
\alpha''=-\varepsilon'_2\varepsilon'_3\,
u'\bigl(
\underbrace{\alpha_0\deleted{\unsure{\tau}}\sigma}
_{\inserted{\alpha\tau\sigma=\alpha\tau'}}
\bigr)\alpha'\underbrace{\sigma\tau}_{\inserted{-\tau'}}
\]
i.e.
\[
(10)\qquad
\alpha''=\varepsilon'_2\varepsilon'_3\,u'(\alpha\tau')
(\alpha'\tau')
\quad\big|\quad\text{Case }\varepsilon_3=-1
\]
a formula which makes sense, since $\alpha\tau'\in\mathcal{G}^{+}$ hence
$u'(\alpha\tau')$ is defined, and the preceding argument proves that
the element in question \deleted{\unsure{is in}}\inserted{satisfies}
\[
(11)\qquad
\alpha''\in\mathcal{G}^{\unsure{\natural|\sharp}}_\tau,\quad
\varepsilon_{\mathcal{G}}(\alpha'')
=\underbrace{\varepsilon_{\mathcal{G}}(\alpha)}_{-1}
\varepsilon_{\mathcal{G}}(\alpha')
\]
by distinguishing cases \unsure{(*)} and $(**)$ of (9).
Same conclusion in the first case where
$\varepsilon_{\mathcal{G}}(\alpha)=1$ (7).
\inserted{This completes the proof.}

\mspage{654}

\scanpage{656}

\uline{Remarks} \deleted{a) For condition b) to be
\unsure{multiplicative}, we do not need condition a) in
\unsure{all} its strength, it suffices that $u$ in
$\mathcal{G}^{+}/\mathcal{G}^{\unsure{\natural|\sharp}}$
normalize the image of $L_\sigma$, i.e. that we have
\[
u(\sigma)\equiv\pm\sigma
\quad(\mathcal{G}^{\unsure{\natural|\sharp}}).
\]}

\deleted{b)} If $\mathcal{G}$ is ``big enough'', the projection homomorphism
is injective

(12)\quad
\diagram{\begin{tikzcd}[ampersand replacement=\&, column sep=large]
Z_{\rho,\sigma}\arrow[r,hook]\&\operatorname{Aut}(\mathcal{G}^{+})
\end{tikzcd}}
\quad (cf. page 578)

\uline{Examples} in the universal case

a) \quad $\mathcal{G}^{\unsure{\natural|\sharp}}=\mathcal{G}^{+}
=\deleted{\unsure{\mathcal{G}}}\mathrm{Sl}(2,\mathbb{Z})^{\wedge}$.
\deleted{Here we have the} \unsure{Condition|condition} (4) d) is empty,
conditions a), b) are equivalent respectively to the fact that
$u(\rho)$ conjugate to $\rho^{\varepsilon_3}$,
$u(\sigma)$ conjugate to $\sigma^{\varepsilon_2}$,
i.e. \unsure{comes back} to definition (6) page 578, which gives
\[
\deleted{Z_{\rho,\sigma}=M(\unsure{\omega|0}).L_0^{\mathcal{G}}
=M^{\mathcal{G}}[0].}
\]
\[
(13)\qquad
Z_{\rho,\sigma}=M(0).L_0^{\mathcal{G}}
=M^{\mathcal{G}}[0]=M[0]
\]

\margin{\uline{N.B.} A priori we rather have
$Z_{\rho,\sigma}\subset\unsure{\widetilde{M}_{0,3}|M_{0,3}}$, but
\diagram{\begin{tikzcd}[ampersand replacement=\&, column sep=large,row sep=large]
M[0]\arrow[r,"\sim"]\arrow[d,equal]\&
\widetilde{M}_{0,3}[0]\arrow[d,equal]\\
M(\unsure{\omega|0}).L_0^{\mathcal{G}}\arrow[r,"\sim"]\&
\unsure{M_{0,3}(0)}.
L_0^{\tikz[baseline=(g.base)]{\node[draw,ellipse,inner sep=1pt]
(g){$\scriptstyle\unsure{\mathcal{G}_{0,3}}$};}}
\end{tikzcd}}}

b) \quad
$\mathcal{G}^{\unsure{\natural|\sharp}}=\Pi=\Pi_0\times\mu$.
We now find
\[
(14)\qquad
Z_{\rho,\sigma}=M(0).L_0^\Pi
=M^\Pi[0]=\unsure{M^{!}[0]|M^{\prime!}[0]}
\]
(denoted $Z_{\rho,\sigma}^{!}$ in $\mathrm{V}$ number$^\circ$ 8).

\margin{In this case, it is immediate that $Z_{\rho,\sigma}$
operates \unsure{trivial\textsuperscript{ly}} even on
$\mathcal{G}^{+}/\mathcal{G}^{\unsure{\natural|\sharp}}_\tau$,
for this is so for the
\deleted{\unsure{N.B.} \unsure{all the subgroups}
\unsure{contained by} a), b) of (4),
$\unsure{\mathcal{G}^{\natural}=M(0)_2.L_0^\Pi|
\mathcal{G}^{\natural}=M(\omega)_2.L_0^\Pi}$}.}
\margin{$\unsure{M^{!}[0]|M^{\prime!}[0]}=M^\Pi[0]$.}
\margin{make explicit again why
\unsure{(i.e. why)} that
$(\varepsilon_0,1,1,1)\notin Z_{\rho,\sigma}$}

c) \quad $\mathcal{G}^{\unsure{\natural|\sharp}}=\Pi_0$.
Conditions a) and b) \inserted{of (4)} have not changed
\unsure{with respect to|with respect to the}
\deleted{\unsure{before}}\inserted{\uline{preceding case}},
since $\Pi=\Pi_0\times\mu$ and $\mu$ is central, on the other hand the

\mspage{655}

\scanpage{657}
\mspage{631}

condition d) becomes distinctly stricter, we find for example that
$Z_{\rho,\sigma}\cap\mathcal{G}^{+}=L_0^{\Pi_0}$
(of index 2 in the corresponding group $L_0^\Pi$ of the preceding case),
and since the operation of $M(0)$ \inserted{$\subset M^!$} on
$\mathcal{G}^{+}/\Pi_0$ is through the character
$\varepsilon_2^{\unsure{*|\#}}$, we see that we now have
\[
(15)\qquad
Z_{\rho,\sigma}=M''(0).L_0^{\Pi_0}=M''[0]_0
\]
\[
\left\{
\begin{array}{l}
M''(0)=\operatorname{Ker}(M(0)\xrightarrow{\varepsilon_2}\mu)\\
M''[0]_0=\operatorname{Ker}(M[0]_0\xrightarrow{\varepsilon_2}\mu)\\
M[0]_0=M(\unsure{\omega|0}).L_0^{\Pi_0}
\end{array}
\right.
\]

\margin{\deleted{This is so if in definition (4) of
$Z_{\rho,\sigma}$ we require in d) that $u$
\unsure{induce|produce} the trivial operation on
$\mathcal{G}^{+}/\mathcal{G}^{\unsure{\natural|\sharp}}$,
a condition that I ended up abandoning as useless.
\unsure{In the new version} the cases
$\mathcal{G}^{\unsure{\natural|\sharp}}=\Pi$ and
$\mathcal{G}^{\unsure{\natural|\sharp}}=\Pi_0$
give the same $Z_{\rho,\sigma}$ (15).
But I would like to obtain
\[
M^{\Pi_0}[0]=\unsure{M[0]^!_0|M[0]_0}
=M(\unsure{\omega|0}).\Pi_0
\]
--- the group of b) is \unsure{too} large, that of c)
\unsure{a little too}
\uline{\unsure{(old version)}}
\inserted{a little too small!}}}

\margin{\deleted{But with the new definition of
$Z_{\rho,\sigma}$ (condition added to (4)) we will finally have
\[
Z_{\rho,\sigma}=M^![0]
\]
as in the preceding case b).
\inserted{We could have \unsure{introduced}
an $\unsure{M^![0]_0|M^{\prime!}[0]_0}
=M(\unsure{0|\omega}).\unsure{L_0^{\Pi_0}|\Pi_0}$.}}}

\margin{This remark about a
$\mathcal{G}^{\unsure{\natural|\sharp}}_0$
now becomes useless.}

These two \inserted{latter} examples suggest the advisability,
in the \inserted{general} case where we would have a subgroup
\deleted{\unsure{$Z$}} $\mathcal{G}^{\unsure{\natural|\sharp}}_0$
of $\mathcal{G}^{\unsure{\natural|\sharp}}$,
\inserted{invariant in $\mathcal{G}$}, such that we have
\[
(16)\qquad
\mathcal{G}^{\unsure{\natural|\sharp}}
=\mathcal{G}^{\unsure{\natural|\sharp}}_0\times\mu
\quad\deleted{\unsure{\subset|\Leftarrow}}
\]
\inserted{(here $\mu$ is the universal central subgroup
$\mu=\{1,\rho^3\}=\{1,\sigma^2\}$)}
of defining $Z_{\rho,\sigma}$ in terms of
$\mathcal{G}^{\unsure{\natural|\sharp}}$, and of
\uline{choosing} the $\alpha,\beta$ of definition (4) in
\deleted{$\unsure{\mathcal{G}^{\natural}_{0\tau}}$}
\[
(17)\qquad
\mathcal{G}^{\unsure{\natural|\sharp}}_{0\tau}
\overset{\mathrm{dfn}}{=}
\mu_\tau.\mathcal{G}^{\unsure{\natural|\sharp}}_0.
\]

\uline{Remark} We are tempted to replace definition (4)
\inserted{(for a) and b))}
by the apparently simpler definition of page 578,
which makes it evident that we have a subgroup.
But they \deleted{\unsure{are compatible}}
\inserted{do not give, even if
$\mathcal{G}^{\unsure{\natural|\sharp}}\subset\Pi$
(\uline{universal case})}
\deleted{with \unsure{if we suppose}}
a trivial action of
\deleted{$\mathcal{G}^{+}/\mathcal{G}^{\unsure{\natural|\sharp}}
\longrightarrow\mathcal{G}_3$}
\inserted{$Z$ on $\mathcal{G}^{+}/\Pi$,
as we \deleted{\unsure{see}}
\inserted{see} immediately,
\deleted{\unsure{$\rho,\sigma$}}}

\mspage{656}

\scanpage{658}

except \unsure{by} \unsure{restricting} ourselves to some, $(u,\mu,\varepsilon_2,\varepsilon_3)$ with $\varepsilon_3=1$,
since $\dot{\rho}^{-1}$ is not equal to $\dot{\rho}$ ! It is true \deleted{that}
(in the universal case for example) that $u(\rho)$ is conjugate \uline{in $\mathcal{G}^{+}$}
to $\rho^{\varepsilon_3}$, but in general it will not be so by an element (say) of $\Pi$.

We set
\[
(18)\qquad
L_0^{\unsure{\natural|\sharp}}
=\deleted{\unsure{L_0}\cap\mathcal{G}^{\unsure{\natural|\sharp}}}
\inserted{\bigl\{l\in L_0^{\mathcal{G}}\mid
(\operatorname{int}(l),1,1,1)\in Z_{\rho,\sigma}\bigr\}}
\qquad
\left(\text{where }L_0^{\mathcal{G}}\subset\mathcal{G}^{+}
\text{ is the \unsure{set} of the }\varepsilon_0^n,\ n\in\widehat{\mathbb{Z}}\right)
\]
\margin{\uline{N.B.} We have
$L_0^{\unsure{\natural|\sharp}}\supset L_0\cap\mathcal{G}^{\unsure{\natural|\sharp}}$,
and the inclusion can be strict, cf. example c) above\ldots
where \uline{$L_0^{\unsure{\natural|\sharp}}=L_0^\Pi$},
$L_0\cap\mathcal{G}^{\unsure{\natural|\sharp}}=L_0^{\Pi_0}$.
\unsure{which} \unsure{is surprising}.}

and we have a canonical homomorphism
\diagram{
\begin{tikzcd}[ampersand replacement=\&]
(19)\quad L_0^{\unsure{\natural|\sharp}}
  \arrow[r,"\unsure{i}"] \& Z_{\rho,\sigma},\\
l \arrow[r,mapsto] \& (\operatorname{int}_{\mathcal{G}^{+}}(l),1,1,1)
\end{tikzcd}
}
injective iff we have
\[
(20)\qquad
L_0^{\unsure{\natural|\sharp}}\cap\operatorname{Centr}(\mathcal{G}^{+})=\{1\}
\]
the image invariant. We set
\[
(21)\qquad
\gamma_{\rho,\sigma}=Z_{\rho,\sigma}/L_0^{\unsure{\natural|\sharp}},
\]
whence
\diagram{
\begin{tikzcd}[ampersand replacement=\&]
(22)\quad 1 \arrow[r] \&
L_0^{\unsure{\natural|\sharp}} \arrow[r] \&
Z_{\rho,\sigma} \arrow[r,"\delta_Z"] \&
\gamma_{\rho,\sigma} \arrow[r] \& 1,
\end{tikzcd}
}
\diagram{
\begin{tikzcd}[ampersand replacement=\&]
 \& \& \mu\\
(23)\quad Z_{\rho,\sigma} \arrow[r] \&
\gamma_{\rho,\sigma}
  \arrow[ur,"\varepsilon_{3\gamma}"]
  \arrow[r,"\chi_\gamma"]
  \arrow[dr,"\varepsilon_{2\gamma}"'] \&
\widehat{\mathbb{Z}}^{*}\\
 \& \& \mu
\end{tikzcd}
}
\unsure{)} \unsure{deduced} from the
\diagram{
\begin{tikzcd}[ampersand replacement=\&, row sep=0pt]
 \& \varepsilon_3\\
(u,\mu,\varepsilon_2,\varepsilon_3)
  \arrow[ur,mapsto] \arrow[r,mapsto] \arrow[dr,mapsto]
  \& \mu\\
 \& \varepsilon_2.
\end{tikzcd}
}

\deleted{
\uline{Remarks} We \deleted{denote} set,
\[
(24)\qquad
\left\{
\begin{array}{lll}
\gamma'_{\rho,\sigma}=\operatorname{Ker}\varepsilon_{3\gamma},&
\gamma''_{\rho,\sigma}=\operatorname{Ker}\unsure{\varepsilon_{2,\sigma}|\varepsilon_{2\gamma}},&
\gamma'''_{\rho,\sigma}=\operatorname{Ker}\varepsilon_{2\gamma}\varepsilon_{3\gamma}\\
Z'_{\rho,\sigma}=\ldots&
Z''_{\rho,\sigma}=\ldots&
Z'''_{\rho,\sigma}=\ldots
\end{array}
\right.
\qquad
(\text{inverse images in }Z_{\rho,\sigma}\unsure{)}
\]
}
\mspage{657}